\documentclass[11pt]{article}

\usepackage{amsmath,amssymb,amstext,amsthm}
\usepackage{geometry}
\usepackage{fancyhdr}
\usepackage[pdftex]{graphicx}
\usepackage{xcolor}

\geometry{
  verbose,
  dvips,
  width=430pt, marginparsep=0pt, marginparwidth=0pt,
  top=90pt, headheight=12pt, headsep=20pt, footskip=30pt, bottom=80pt}

\setlength{\parskip}{\medskipamount}
\setlength{\parindent}{0pt}

\linespread{1.05}

\newtheorem{theorem}{Theorem}
\newtheorem{lemma}[theorem]{Lemma}
\newtheorem{proposition}[theorem]{Proposition}
\newtheorem{corollary}[theorem]{Corollary}
\newtheorem{conjecture}[theorem]{Conjecture}
\newtheorem{obs}{Observation}
\newtheorem{clm}{Claim}
\newtheorem{ques}{Question}
\newtheorem{problem}{Problem}

\newtheorem{ex}{Example}


\newcommand{\spc}{\hspace{0.08in}}
\newcommand{\vs}{\vspace*{.1in}}
\newcommand{\E}{\mathbb{E}}
\newcommand{\Prob}{\mathbb{P}}
\newcommand{\edit}[1]{\emph{\textcolor{blue}{#1}}}


\renewenvironment{proof}{{\noindent \textbf{\textit{Proof.}}}}
{\hfill $\Box$ \vspace*{0.1in}}
\newenvironment{proofc}{{\noindent \textit{Proof of Claim.}}}
{\hfill $\Box$}


\begin{document}

\begin{LARGE}\begin{center}\bf 12.4 Cross Product \end{center}
\end{LARGE}

\vs\vs



\section{Warm Up}

\textbf{Question}: When in calculus 1 or 2 did we care about being ``perpendicular"? Why may we need this in 3-D?

\textbf{Problem 1}: Find the determinant of 
$\begin{bmatrix}
1 & -2 \\
4 & 3
\end{bmatrix}$


\section{Motivation}
Suppose I have two vectors $u$ and $v$. It is seemingly useful to know a vector $w$ which is perpendicular to \emph{BOTH} $u$ and $v$. From our previous chapter we know that in order for this to happen we would need $u \cdot w = 0$ and $v \cdot w = 0$. This resulting vector we will call the \emph{cross product} of $u$ and $v$.



\section{Definitions \& Theorems}

\begin{enumerate}
    \item If $a = \langle a_1,a_2,a_3 \rangle$ and $b = \langle b_1,b_2,b_3 \rangle$, then the cross product of $a$ and $b$ is the vector 

    \begin{equation*}
        a \times b = \langle a_2b_3-a_3b_2, a_3b_1-a_1b_3, a_1b_2 - a_2b_1 \rangle.
    \end{equation*}


    \item The cross product can be found by taking the determinant of the 3 by 3 matrix whose first row is $(i,j,k)$, second row is $a = (a_1,a_2,a_3)$ and third row is $b = (b_1,b_2,b_3)$.

    \item  We can also leverage the cross product to find the angle $\theta$ between two vectors. We get the following theorem.

    \begin{theorem}
        If $\theta$ is the angle between vectors $a$ and $b$, then $|a\times b| = |a||b|\sin(\theta)$.
    \end{theorem}

    \begin{proof}
    \begin{align*}
        |a\times b|^2 &= (a_2b_3-a_3b_2)^2+(a_3b_1-a_1b_3)^2+(a_1b_2-a_2b_1)^2 \\
        &= a_2^2b_3^3-2a_2a_3b_2b_3+a_3^2b_2^2+a_3^2b_1^2-2a_1a_3b_1b_3+a_1^2b_3^2+a_1^2b_2^2 - 2a_1a_2b_1b_2 +a_2^2b_1^2 \\
        &=(a_1^2+a_2^2+a_3^2)(b_1^2+b_2^2+b_3^2)-(a_1b_1+a_2b_2+a_3b_3)^2\\
        &=|a|^2|b|^2-(a \cdot b)^2 \\
        &= |a|^2|b|^2-|a|^2|b|^2\cos^2(\theta)\\
        &=|a|^2|b|^2(1-\cos^2(\theta))\\
        &=|a|^2|b|^2\sin^2(\theta)
    \end{align*}
    \end{proof}

    \item Since the angle between parallel vectors is 0 and $\sin(0)=0$, we see that two vectors are parallel if and only if $a\times b = 0$.

    \item The geometric representation of the cross product gives us the following insight. The magnitude of the cross product $|a \times b|$ is equal to the area of the parallelogram determined by $a$ and $b$. This can easily seen by taking the sin of the angle in our parallelogram and solving for base times height.

    \item While the proof of this next theorem is not super enlightening, it is interesting to see how we can combine these notions of products and find the volume of a 3-D parallelogram.

    \begin{theorem}
        The volume of the parallelepiped determined by the vectors $a,b$, and $c$ is $V=|a \cdot (b\times c)|$.
    \end{theorem}

    \item  Torque $\tau = r \times F$ where $r$ is a position vector and $F$ is a force vector.
\end{enumerate}



\section{Examples}


\begin{problem}
    Find the determinant of the following $3\times 3$ matrix.

    \begin{equation*}
        \begin{bmatrix}
            1&-2&2 \\
            3 &1&2 \\
            -2 & -2 &1 
        \end{bmatrix}
    \end{equation*}
\end{problem}

\vspace{3cm}


\begin{problem}
    Find $\langle 5,1,-3 \rangle \times \langle 4,-4,2 \rangle$.
\end{problem}

\vspace{6cm}
\begin{problem}
Find a vector which is perpendicular to the plane which passes through the points $A(1,4,6)$, $B(-2,5,-1)$, and $C(1,-1,1)$.
\end{problem}

\vspace{5cm}


\begin{problem}
    Find the area of the triangle whose two sides are formed by $u = 4i-2k$ and $v=-i+3j+2k$.
\end{problem}

\vspace{5cm}

\begin{problem}
    Find the volume of the paralellepiped determined by the vectors $a= \langle 2,0,-3 \rangle$, $b=\langle 5,1,1 \rangle$, and $c=\langle -2,3,1 \rangle$.
\end{problem}


\vspace{8 cm}

\begin{problem}
    A bolt is tightened by applying a 40N force to a .25-m wrench. Find the magnitude of the torque at the center of the bolt.
\end{problem}

\newpage

\section{Scratch Work}


\end{document} 